% ============================================================
% Capítulo 20: Astronomía Multi-Mensajero y Fronteras
% Relatividad, Cosmología y Ondas Gravitacionales
% Daniel Soto Villada — Universidad de Antioquia
% ============================================================

\chapter{Astronomía Multi-Mensajero y Fronteras de la Investigación}
\label{cap:multimensajero_fronteras}

La detección de GW170817 junto con GRB 170817A y su kilonova asociada transformó la astronomía de ondas gravitacionales en una ciencia genuinamente multi-mensajero \cite{Abbott2017BNS, Abbott2017multimessenger, Abbott2017H0}. Desde entonces, las preguntas centrales no son solo ``¿hay señal?'', sino ``¿qué física nueva podemos restringir al combinar mensajeros?'', ``¿cómo escalar precisión con redes globales?'' y ``¿qué ventanas abrirán LISA y PTA?'' \cite{LISA2017, NANOGrav2023, Maggiore2018}.

Este capítulo sintetiza el estado del arte y traza líneas de investigación para la próxima década en gravitación, astrofísica compacta y cosmología de precisión.

% ===========================================================
\section{De la detección a la astronomía integrada}
\label{sec:deteccion_astronomia_integrada}
% ===========================================================

\begin{definicion}{Astronomía multi-mensajero}{def:multi_mensajero}
Es la inferencia física conjunta a partir de señales gravitacionales, electromagnéticas, neutrinos y/o rayos cósmicos, explotando complementariedad observacional para romper degeneraciones.
\end{definicion}

\begin{observacion}{Ganancia epistemológica}{obs:ganancia_epistemologica}
Cada mensajero mide subconjuntos distintos de parámetros; la combinación produce restricciones más robustas que cualquier canal por separado.
\end{observacion}

Ejemplos típicos de complementariedad:
\begin{itemize}[leftmargin=2em]
  \item GW: distancia luminosa y dinámica orbital,
  \item EM: redshift, ambiente anfitrión, física radiativa,
  \item neutrinos: cronología y microfísica de colapso.
\end{itemize}

% ===========================================================
\section{GW170817 como punto de inflexión}
\label{sec:gw170817_inflexion}
% ===========================================================

\begin{definicion}{Evento ancla}{def:evento_ancla}
Un evento ancla es una detección con riqueza observacional suficiente para calibrar metodologías en múltiples dominios (detección, localización, inferencia, seguimiento EM, cosmología).
\end{definicion}

GW170817 cumple este rol por:
\begin{enumerate}[leftmargin=2em, label=\textbf{\arabic*.}]
  \item detección GW robusta de una BNS,
  \item contraparte gamma en segundos,
  \item evolución kilonova en óptico/infrarrojo,
  \item identificación de galaxia anfitriona y redshift.
\end{enumerate}

\begin{importante}
La combinación GW+EM permitió medir distancia y redshift en un mismo evento sin usar la escalera de distancia tradicional, validando operacionalmente el método de sirenas estándar.
\end{importante}

% ===========================================================
\section{Restricciones a propagación y velocidad de ondas gravitacionales}
\label{sec:restricciones_propagacion}
% ===========================================================

\begin{definicion}{Retardo inter-mensajero}{def:retardo_intermensajero}
Es la diferencia de tiempo de llegada entre señales de distintos mensajeros emitidos en un mismo fenómeno astrofísico, modelada como suma de retardo intrínseco de emisión y posible retardo de propagación.
\end{definicion}

Para distancia $D$ y velocidades $v_{\mathrm{GW}},c$:
\begin{equation}
\Delta t_{\mathrm{prop}} \approx D\left(\frac{1}{v_{\mathrm{GW}}}-\frac{1}{c}\right).
\end{equation}

El pequeño desfase observado en GW170817/GRB170817A impuso límites severos a desviaciones de $v_{\mathrm{GW}}$ respecto a $c$, restringiendo clases de teorías modificadas de gravedad.

\begin{observacion}{Dependencia de modelado astrofísico}{obs:dependencia_modelado}
La interpretación cuantitativa requiere separar cuidadosamente retardo de emisión en el jet y retardo de propagación fundamental.
\end{observacion}

% ===========================================================
\section{Materia densa y ecuación de estado en estrellas de neutrones}
\label{sec:eos_materia_densa}
% ===========================================================

\begin{definicion}{Deformabilidad mareal efectiva}{def:lambda_tidal}
La deformabilidad mareal codifica cuánto se deforma una estrella de neutrones en el campo gravitacional de su compañera y deja huella en la fase de inspiral tardío.
\end{definicion}

Las restricciones combinadas GW+EM informan:
\begin{itemize}[leftmargin=2em]
  \item radios típicos de estrellas de neutrones,
  \item rigidez de la ecuación de estado,
  \item masa máxima estable frente a colapso.
\end{itemize}

\begin{proposicion}{Sinergia observacional}{prop:sinergia_eos}
Datos GW acotan deformabilidad macroscópica; datos EM de kilonova y remanente aportan información sobre eyección y composición, reforzando inferencia sobre microfísica nuclear.
\end{proposicion}

% ===========================================================
\section{Cosmología con sirenas estándar}
\label{sec:cosmologia_sirenas_estandar}
% ===========================================================

\begin{definicion}{Sirena estándar brillante y oscura}{def:sirenas_brillante_oscura}
\textbf{Brillante:} evento con contraparte y anfitriona identificada (redshift directo). \\
\textbf{Oscura:} evento sin contraparte inequívoca; el redshift se infiere estadísticamente con catálogos de galaxias.
\end{definicion}

Para eventos de bajo redshift:
\begin{equation}
d_L \simeq \frac{c\,z}{H_0},
\end{equation}
por lo que la medida de $d_L$ (GW) y $z$ (EM/catálogo) permite inferir $H_0$.

\begin{observacion}{Escalamiento de precisión}{obs:escalamiento_h0}
La precisión en $H_0$ mejora aproximadamente como $1/\sqrt{N_{\mathrm{eff}}}$ cuando se combinan eventos independientes con sistemáticas controladas.
\end{observacion}

\subsection{Retos actuales}

\begin{itemize}[leftmargin=2em]
  \item degeneración distancia--inclinación en evento individual,
  \item incompletitud y sesgo de catálogos de galaxias para sirenas oscuras,
  \item peculiar velocities en régimen de bajo redshift.
\end{itemize}

% ===========================================================
\section{Mapeando el cielo dinámico: localización y respuesta rápida}
\label{sec:localizacion_respuesta_rapida}
% ===========================================================

La ciencia multi-mensajero exige tiempos de respuesta de minutos a horas:
\begin{itemize}[leftmargin=2em]
  \item alertas de baja latencia,
  \item mapas probabilísticos de localización,
  \item priorización de telescopios por área, profundidad y ventana temporal.
\end{itemize}

\begin{definicion}{Cadena de seguimiento}{def:cadena_seguimiento}
Secuencia coordinada de publicación de alerta, refinamiento de localización, estrategia de apuntado y actualización de clasificación astrofísica en tiempo casi real.
\end{definicion}

\begin{importante}
La calidad de la localización no depende solo de SNR, sino de geometría de red, disponibilidad de detectores y calibración en el momento del evento.
\end{importante}

% ===========================================================
\section{Nuevas bandas: LISA y arreglos de púlsares (PTA)}
\label{sec:lisa_pta}
% ===========================================================

\subsection{LISA y la banda milihertz}

LISA observará binarias masivas de agujeros negros, sistemas galácticos compactos y potencialmente fondos cosmológicos en mHz \cite{LISA2017}.

\begin{observacion}{Astronomía multibanda}{obs:multibanda}
Algunas fuentes podrán seguirse desde mHz (LISA) hasta banda terrestre, permitiendo predicción previa de coalescencia y campañas EM planificadas.
\end{observacion}

\subsection{PTA y la banda nanohertz}

PTA explota la estabilidad temporal de púlsares milisegundo para detectar correlaciones angulares inducidas por ondas gravitacionales de ultra baja frecuencia \cite{NANOGrav2023}.

\begin{definicion}{Firma de Hellings--Downs}{def:hellings_downs}
Patrón angular esperado de correlación temporal entre pares de púlsares en presencia de un fondo estocástico gravitacional isotrópico.
\end{definicion}

% ===========================================================
\section{Fronteras teóricas y observacionales}
\label{sec:fronteras_teoricas_observacionales}
% ===========================================================

Líneas activas de investigación:
\begin{itemize}[leftmargin=2em]
  \item tests de relatividad general en régimen dinámico fuerte,
  \item búsqueda de polarizaciones extra y disipación anómala,
  \item demografía de agujeros negros de masa intermedia y supermasiva,
  \item fondos cosmológicos del universo temprano,
  \item inferencia cosmológica conjunta con CMB/BAO/supernovas.
\end{itemize}

\begin{observacion}{Desafío metodológico transversal}{obs:desafio_metodologico}
El límite actual suele ser inferencial, no detectivo: controlar sistemáticas, sesgos de selección y consistencia entre instrumentos es tan crítico como mejorar sensibilidad bruta.
\end{observacion}

% ===========================================================
\section{Ejemplo desarrollado: estrategia multi-mensajero para una BNS cercana}
\label{sec:ejemplo_estrategia_multimensajero}
% ===========================================================

\begin{ejemplo}{Plan científico en las primeras 48 horas}{ej:plan_48h}
Considere una BNS detectada por una red de tres interferómetros con localización inicial de 120 deg$^2$ y distancia $d_L$ moderadamente bien medida.

\textbf{Paso 1: priorización inicial (0--2 h).} \\
Cruzar mapa GW con catálogos de galaxias y seleccionar campos de mayor probabilidad ponderando visibilidad y magnitud límite.

\textbf{Paso 2: búsqueda transitoria (2--24 h).} \\
Ejecutar barridos ópticos/IR para identificar fuentes nuevas con evolución compatible con kilonova.

\textbf{Paso 3: confirmación y espectros (24--48 h).} \\
Obtener espectros y fotometría multibanda para separar supernovas comunes de candidatos kilonova.

\textbf{Paso 4: retorno a inferencia conjunta.} \\
Incorporar redshift anfitrión y datos EM en análisis bayesiano para refinar inclinación, distancia y parámetros cosmológicos.
\end{ejemplo}

% ===========================================================
\section{Síntesis final}
\label{sec:sintesis_cap20}
% ===========================================================

Conclusiones del capítulo:
\begin{enumerate}[leftmargin=2em, label=\textbf{\arabic*.}]
  \item La astronomía multi-mensajero transforma detecciones GW en física integrada de alta precisión.
  \item GW170817 estableció un estándar metodológico para pruebas de propagación, materia densa y cosmología.
  \item Sirenas estándar, brillantes y oscuras, son un pilar emergente para medir parámetros cosmológicos.
  \item LISA y PTA abrirán ventanas de frecuencia complementarias con fuerte potencial de descubrimiento.
  \item El principal reto de la próxima década es controlar sistemáticas y coordinar redes observacionales globales.
\end{enumerate}

Con este capítulo concluye el recorrido del libro: de los fundamentos geométricos y relativistas hasta las fronteras experimentales y cosmológicas de la gravitación moderna.

% ===========================================================
\section*{Problemas propuestos}
\addcontentsline{toc}{section}{Problemas propuestos}
% ===========================================================

\begin{enumerate}[leftmargin=2.2em, label=\textbf{P\arabic*.}]

\item \textbf{Complementariedad de mensajeros.} \\
Dé un ejemplo concreto de parámetro físico que sea débilmente constriñido por GW solo y mejore al combinar con EM.

\item \textbf{Velocidad de ondas gravitacionales.} \\
Explique cualitativamente cómo un retardo GW--gamma pequeño restringe teorías de gravedad modificada.

\item \textbf{Sirenas estándar oscuras.} \\
Discuta el papel de la completitud de catálogos galácticos en la inferencia de $H_0$ sin contraparte inequívoca.

\item \textbf{Materia densa.} \\
Relacione deformabilidad mareal con información sobre rigidez de la ecuación de estado.

\item \textbf{Respuesta rápida.} \\
Proponga tres criterios para priorizar campos de observación EM tras una alerta GW.

\item \textbf{Multibanda LISA--terrestre.} \\
Describa una ventaja científica de detectar una misma fuente en mHz y luego en decenas de Hz.

\item \textbf{PTA y fondos estocásticos.} \\
¿Por qué la correlación angular entre púlsares es central para distinguir señal común de ruido intrínseco?

\end{enumerate}
